\documentclass[11pt]{article}
\usepackage[a4paper,margin=25mm]{geometry}
\usepackage[T1]{fontenc}
\usepackage[utf8]{inputenc}
\usepackage{lmodern,microtype}
\usepackage{amsmath,amssymb,amsthm,mathtools}
\usepackage{booktabs,array,longtable,enumitem,float}
\usepackage{xcolor,xurl}
\usepackage[hidelinks]{hyperref}
\usepackage{fancyhdr}
\usepackage{listings}
\definecolor{ink}{HTML}{193A54}
\hypersetup{pdftitle={Power-Conjugacy Saturation for Unknot Recognition},pdfauthor={Research contribution prepared with ChatGPT for the ProveIt project}}
\pagestyle{fancy}\fancyhf{}\fancyhead[L]{\small Power-conjugacy saturation}\fancyhead[R]{\small ProveIt research continuation}\fancyfoot[C]{\thepage}
\setlength{\headheight}{14pt}
\setlength{\parskip}{3pt}
\setlist{nosep,leftmargin=*}
\lstset{basicstyle=\ttfamily\small,columns=fullflexible,breaklines=true,frame=single,framerule=0.3pt,showstringspaces=false}
\newtheorem{theorem}{Theorem}[section]
\newtheorem{lemma}[theorem]{Lemma}
\newtheorem{proposition}[theorem]{Proposition}
\newtheorem{corollary}[theorem]{Corollary}
\theoremstyle{definition}\newtheorem{definition}[theorem]{Definition}
\theoremstyle{remark}\newtheorem{remark}[theorem]{Remark}
\newcommand{\Z}{\mathbb Z}\newcommand{\Q}{\mathbb Q}\newcommand{\F}{\mathbb F}
\newcommand{\bits}{\operatorname{bits}}\newcommand{\poly}{\operatorname{poly}}
\newcommand{\Kh}{\widetilde{\mathrm{Kh}}}
\newcommand{\code}[1]{\texttt{#1}}
\newcommand{\G}{\mathcal G}
\newcommand{\dead}{\mathcal D}
\newcommand{\M}{\mathsf M}
\newcommand{\TestCount}{24}
\newcommand{\AuditCount}{3227}

\title{\color{ink}Power-Conjugacy Saturation\\for Unknot Recognition\\[5pt]\large Sharp graph certificates, polynomial deletion closure,\\and binary power-pair kernels}
\author{Research contribution prepared with ChatGPT\\for Vladimir Reshetnikov's \textit{ProveIt} project}
\date{8 October 2026\\\small Computational records: 9 October 2026 UTC}
\begin{document}
\maketitle
\begin{abstract}
We develop a source-checkable group-presentation operation complementary to
ProveIt's primitive-pair and unit-coordinate forest contractions. Verified
relations between conjugates of generator powers define a labelled graph.
Using the established balancedness property of classical knot groups, an
inconsistent absolute exponent ratio around a cycle forces every generator
in its connected component to be trivial. Plain power equalities admit an
additional signed-cycle test. We prove that these tests are complete for the
graph abstraction: every surviving component has a simultaneous nontrivial
realization inside the trefoil group when the discarded conjugator labels
are treated as independent ambient elements. Thus the distinction between
absolute and signed consistency is necessary, not an implementation choice.

We give short independently replayable certificates, deterministic exact
fallback behind a fixed-prime sieve, and a conservative polynomial bit bound
for complete deletion-only saturation of literal presentations. Binary
power-pair families permit polynomial-cost elimination where the existing
local primitive-forest eligibility conditions provide no move; a gcd-of-minors
criterion also yields an unconditional two-generator triviality theorem.
A self-contained Python package includes an independent source replayer,
a small reduced-Khovanov reference oracle, \TestCount{} passing test methods,
and a \AuditCount{}-input braid audit. The audit has no incorrect positive
verdict, but adds no coverage over pure-power deletion on that corpus.
Controlled sieve measurements retain both gains and arithmetic aliasing
failures. These are local and restricted-input results, not an unrestricted
quasi-polynomial unknot-recognition theorem or a measured speedup of the
maintained recognizer.
\end{abstract}
\clearpage
\begingroup\small\setlength{\parskip}{0pt}
\tableofcontents
\endgroup
\newpage

\section{Research target and relationship to the repository}
\label{sec:target}
The input problem is to decide whether an explicitly encoded classical knot
diagram in $S^3$, with $n$ crossings, represents the unknot. The intended
asymptotic target is $2^{(\log n)^{O(1)}}$. A correct local simplifier is useful
for this target only when its own bit cost, its effect on intermediate
representations, and the cost of discovering the simplifications are all
accounted for. A terminating sequence of presentation changes is not by
itself a running-time theorem.

The repository inspection was pinned to commit
\begin{center}\small\texttt{13d60da8e683ca02b11c6b301c068bda929014e7}.\end{center}
We read the current synthesis introduction and the primitive-projection,
primitive-forest and group-certificate sections, together with the synthesis
README and report 45's README \cite{Repo}. The relevant maintained ideas
are already substantial. A rank-two primitive-power terminal is implemented;
disjoint primitive-pair quotients and overlapping unit-coordinate forests
have independent replay; raw circuits and explicit normalization boundaries
are distinguished; and later synthesis notes include arbitrary acyclic
singleton-definition batches. Repeating the rank-two terminal would not
constitute a continuation.

The present target is the restriction that the forest rule requires acyclic
unit-coordinate dependencies. In a forest edge $c=p^k$, the child can be
replaced by a power of its parent. More general relations such as
\[
 x^a=u y^b u^{-1}
\]
need not solve either generator as an integral power. They may nevertheless
force both generators to be trivial after comparison with other relations.
A cycle can supply the missing information. We study that case without
pretending that every overlapping nonunit relation is a safe forest edge.

\paragraph{Contribution and attribution.}
The knot-group balancedness lemma used below is prior work, as are Artin
presentations, the cyclic-knot-group criterion and Khovanov unknot detection
\cite{HMT,Gemein,Papakyriakopoulos,KM}. The contribution here is the exact
scope of the labelled-graph inference, its sharpness by a fixed-knot
countermodel, its certificate and saturation interfaces, and the associated
restricted complexity results. The gcd-of-minors argument is elementary
and is presented with its full proof, not as an independently established
literature-priority claim. No exhaustive novelty search or external referee
validation is claimed.

\paragraph{The global target remains distinct.}
Lackenby's inspected 2026 preprint estimates the number of passes through a
hierarchy algorithm in terms of a hierarchy-length bound $L$ and a surface
complexity bound $g$. Its Section 9 separates that estimate from the bit
cost of the operations and gives $L(g+1)^L$ in Proposition 9.1
\cite{Lackenby}. The present paper supplies neither the missing global
hierarchy bounds nor a replacement complete quasi-polynomial search. It
isolates a presentation operation whose saturation does not require a
logarithmic-depth hypothesis merely to bound its own repeated execution.

\section{Source hypotheses and the topological input}
\label{sec:background}
Throughout, a \emph{classical knot group} means
$G(K)=\pi_1(S^3\setminus K)$ for a single classical knot $K$. A source-replayed
presentation is a full presentation known, through reconstruction and checked
changes of presentation, to present this group. It is not an arbitrary
presentation accompanied by a Boolean assertion of torsion-freeness.

\begin{lemma}[Established knot-group facts]\label{lem:balanced}
A classical knot group is torsion-free. If $x,u\in G(K)$ and $a,b$ are
nonzero integers with
\[
 x^a=u x^b u^{-1},
\]
then either $x=1$ or $|a|=|b|$.
\end{lemma}
The second statement is Lemma 2.3 of the arXiv version of
Himeno--Motegi--Teragaito \cite{HMT}, with the conjugator inverted to match
our convention. We refer to it as \emph{balancedness}. It is not inferred
from torsion-freeness alone.

For example, consider affine maps of $\Q$ with $x(t)=t+1$ and $u(t)=t/2$.
Then $x=u x^2u^{-1}$ and $x\ne1$. The subgroup generated by these maps is
torsion-free: a nonidentity map with slope $2^k$, $k\ne0$, has infinite
order, and a nonzero translation also has infinite order. Consequently,
a purported verifier that replaces ``classical knot group'' by merely
``torsion-free group'' would be unsound for the absolute-ratio rule.

\begin{proposition}[Cyclic group terminal]\label{prop:terminal}
If a source-replayed presentation of $G(K)$ has one generator and every
remaining relator has exponent sum zero in that generator, then $K$ is
the unknot.
\end{proposition}
\begin{proof}
A word in a single generator is its exponent-sum power. Thus every relator
is the identity in the free cyclic group and the presented group is $\Z$.
The boundary homomorphism $\Z^2\to\pi_1(E(K))=\Z$ is not injective.
The Loop Theorem supplies a compressing disc for the boundary torus. The
knot exterior is irreducible: a sphere in the exterior has a ball on the
side not containing the connected knot. A compressible torus boundary in
this irreducible exterior makes the exterior a solid torus, equivalently
providing the longitude disc and hence a spanning disc for $K$. This is
the standard cyclic-knot-group criterion \cite{Papakyriakopoulos}; the
maintained group checker uses the same full-group criterion \cite{Repo}.
\end{proof}

\paragraph{Two invalid shortcuts.}
First, $G(K)^{\mathrm{ab}}\cong\Z$ holds for every knot, so abelianization
alone cannot authorize Proposition~\ref{prop:terminal}. Second,
torsion-freeness does not imply that $v^d=w^d$ forces $v=w$.
The distinction will be especially important for inverse conjugacy and
for the power-pair arguments below. Those arguments compare powers of the
\emph{same} element or use explicit B\'ezout combinations of equal powers;
they do not assume a general unique-root property.

\section{The power-conjugacy graph and its sound inference rules}
\label{sec:calculus}
\begin{definition}[Verified power-conjugacy edge]
Let $X=\{x_v:v\in V\}$ be current live generators. An oriented edge $e:s\to t$
has nonzero signed integers $(a_e,b_e)$ and a checked relation
\begin{equation}\label{eq:edge}
 x_s^{a_e}=u_e x_t^{b_e}u_e^{-1}.
\end{equation}
It is \emph{plain} only if $u_e=1$ is established. The supplied implementation
uses the literal empty word for this flag. The underlying graph is
undirected for connectivity purposes and permits loops and parallel edges.
\end{definition}
An inverse traversal interchanges $(a_e,b_e)$ and uses $u_e^{-1}$ as its
conjugator. A relator
\begin{equation}\label{eq:donor}
 (x_s^a u x_t^{-b}u^{-1})^d=1,\qquad d>0,
\end{equation}
provides an edge by torsion-freeness. The source replayer must check the
relator, the repeated-root decomposition and all cuts before using this
consequence. The graph routine itself has no access to the knot source.

\begin{lemma}[Path composition]\label{lem:path}
For a directed traversal $P$ with consecutive exponent pairs
$(a_1,b_1),\ldots,(a_k,b_k)$, write
$A_P=\prod_i a_i$, $B_P=\prod_i b_i$, and let $U_P$ be the product of the
traversed conjugators in path order. Then
\begin{equation}\label{eq:path}
 x_{\mathrm{start}(P)}^{A_P}
 =U_P x_{\mathrm{end}(P)}^{B_P}U_P^{-1}.
\end{equation}
\end{lemma}
\begin{proof}
For two successive edges, raise the first equality to $a_2$ and the second
to $b_1$. The identity $(u z u^{-1})^k=u z^k u^{-1}$, valid for every
integer $k$, gives
\[
 x_0^{a_1a_2}
 =u_1 x_1^{b_1a_2}u_1^{-1}
 =u_1u_2 x_2^{b_1b_2}u_2^{-1}u_1^{-1}.
\]
Induction proves the statement. No extraction of roots and no multiplication
of noncommuting base elements is used. In particular, the conjugators are
not exponentiated when paths are composed.
\end{proof}

\begin{lemma}[Propagation of triviality]\label{lem:propagate}
For a verified edge in a torsion-free group, $x_s=1$ if and only if $x_t=1$.
\end{lemma}
\begin{proof}
If $x_s=1$, Equation~\eqref{eq:edge} implies $x_t^{b_e}=1$; since
$b_e\ne0$, torsion-freeness gives $x_t=1$. Reverse the edge for the converse.
\end{proof}

\begin{definition}[Two consistency conditions]\label{def:consistency}
A connected component $C$ is \emph{absolutely consistent} if there are
$q_v\in\Q_{>0}$ satisfying
\begin{equation}\label{eq:abs}
 |a_e|q_s=|b_e|q_t \qquad(e\in E(C)).
\end{equation}
Its plain-edge subgraph is \emph{sign-consistent} if there are
$\epsilon_v\in\{1,-1\}$ satisfying
\begin{equation}\label{eq:sign}
 \epsilon_t=\operatorname{sign}(a_eb_e)\epsilon_s
 \qquad(e\text{ plain}).
\end{equation}
The choices in different connected components of the plain subgraph are
independent. A full component is \emph{bad} if absolute consistency fails
or plain sign consistency fails. Let $\dead$ be the union of its vertices
over all bad components.
\end{definition}

\begin{theorem}[Sound graph annihilation]\label{thm:sound}
For edges verified in a source-replayed classical knot group, every
$x_v$ with $v\in\dead$ is the identity. Adjoining these identities,
deleting the corresponding generators, and substituting the identity for
them in every retained relator preserves the presented knot group.
\end{theorem}
\begin{proof}
Choose a spanning tree of a full component. Assign $q=1$ at its root and
propagate the rational ratios $|a|/|b|$ along tree edges. A failed nontree
edge gives a fundamental cycle $P$ with $|A_P|\ne|B_P|$. Conversely, if
all such edges agree, the propagated values satisfy every edge. Thus
absolute inconsistency has a simple-cycle witness. By
Lemma~\ref{lem:path}, its base generator has conjugate powers with
different absolute exponents and is trivial by Lemma~\ref{lem:balanced}.
Lemma~\ref{lem:propagate} kills the entire component.

If the absolute test succeeds but the plain sign test fails, repeat the
tree argument on a connected component of the plain subgraph. It supplies
a plain cycle with $A_P=-B_P\ne0$. Its conjugator is the identity, so
$x^{A_P-B_P}=1$; torsion-freeness gives $x=1$. Propagation again uses
all edges, not just the plain ones.

Each resulting identity already holds in the full presented group.
Adjoining it therefore changes no quotient. Tietze elimination of the
identity generators means deleting their letters from \emph{all} relators.
Other relators are not discarded simply because they were not selected as
donors. The group remains isomorphic to the same source knot group, so the
hypotheses are retained for the next pass.
\end{proof}

\paragraph{Elementary examples.}
The plain relations $x^2=y^3$ and $x^5=y^7$ give a two-edge cycle with
products $14$ and $15$, so $x=y=1$. The same conclusion holds if the two
equalities have unrelated conjugators, provided the ambient group is a
classical knot group. In contrast, $x=u x^{-1}u^{-1}$ is not a valid
annihilation certificate for arbitrary knot groups. The next section
shows that this last distinction is optimal for the graph abstraction.

\section{Sharpness: every survivor is realizable in one fixed knot group}
\label{sec:sharp}
The graph drops most source information. In particular, it forgets the
words spelling the $u_e$, dependencies among those words, peripheral data,
and the requirement that the named vertices generate the entire group.
Sharpness must be stated for that abstraction, not for the original
presentation.

\begin{definition}[Abstract knot-model]
An abstract knot-model of a labelled graph is a classical knot group $G(K)$,
assigned elements $x_v\in G(K)$, and independently chosen ambient elements
$u_e\in G(K)$ satisfying all edge relations. Plain edges have $u_e=1$.
There is no assertion that the $x_v$ generate $G(K)$ or that a specified
source word evaluates to the chosen $u_e$.
\end{definition}

\begin{theorem}[Exact forced-identity set]\label{thm:sharp}
A vertex is the identity in every abstract knot-model if and only if it
belongs to $\dead$. Moreover, a single model in the trefoil group makes
all vertices outside $\dead$ simultaneously nontrivial.
\end{theorem}
\begin{proof}
The ``if'' assertion is Theorem~\ref{thm:sound}. For the converse use the
fixed group
\[
 T=\langle a,b\mid a^2=b^3\rangle
\]
of the trefoil. Put $g=[a,b]=a^{-1}b^{-1}ab$. The common element $a^2=b^3$
is central, and the commutator identity gives
\[
 1=[a^2,b]=(a^{-1}ga)g.
\]
Hence $a^{-1}ga=g^{-1}$; conjugation by $a$ also inverts $g$ because
$a^2$ is central. The quotient $T\to C_2*C_3$ sends $g$ to the nonempty
reduced alternating word $a b^{-1} a b$, so $g\ne1$. Torsion-freeness of
$T$ makes $g$ infinite-order. This is the trefoil specialization of the
standard even-torus-knot example in \cite{HMT}; the argument is included
to specify exactly what the construction uses.

For each good component, choose positive rational potentials $q_v$ from
Equation~\eqref{eq:abs}, clear a common denominator, and obtain positive
integers $n_v$. Choose signs $\epsilon_v$ satisfying the plain constraints.
Set
\[
 x_v=g^{\epsilon_v n_v}.
\]
Every such element is nontrivial. For an edge let
$A=a_e\epsilon_s n_s$ and $B=b_e\epsilon_t n_t$.
Absolute consistency says $|A|=|B|$. If $A=B$, choose $u_e=1$; if $A=-B$,
choose $u_e=a$. In either case $g^A=u_e g^B u_e^{-1}$. For a plain edge,
Equation~\eqref{eq:sign} ensures $A=B$, so its conjugator is indeed the
identity. On every bad component assign all vertex elements to $1$ and
all conjugators to $1$. Different full components have no edges between
them, completing the simultaneous model.
\end{proof}

\begin{corollary}[A precise algorithmic barrier]\label{cor:barrier}
No sound rule using only the graph's endpoint, exponent and plain-flag
data can universally delete an additional surviving vertex in classical
knot groups. More coverage requires additional source information or a
change of presentation, not a more aggressive interpretation of these
same labels.
\end{corollary}

This is a barrier for the \emph{abstraction}, not a lower bound for unknot
recognition. The trefoil assignment may violate the actual source
conjugator words or its abelianization marking. The original presentation
can therefore contain decisive information that this graph deliberately
forgets.

\paragraph{A topologically conditional strengthening.}
In a knot group having no nontrivial element conjugate to its inverse,
any signed inconsistency around a full graph cycle is sufficient: a
relation $x^A=u x^{-A}u^{-1}$ makes the nonidentity element $x^A$ conjugate
to its inverse. For example, the hyperbolic-knot case has this property
by \cite[Corollary 1.8]{HMT}. The source must establish the stronger
hypothesis; the implementation does not assume hyperbolicity. A generic
``no torsion'' flag cannot replace it.

\section{Algorithms, short certificates and bit cost}
\label{sec:algorithms}
\subsection{Exact discovery without factoring exponents}
The basic algorithm first forms full connected components. Within each,
a breadth-first search propagates positive fractions using $|a|/|b|$.
At an inconsistent edge, the two paths to the tree root and that edge
supply a simple cycle. The whole component is marked for deletion and
its remaining rational queries are skipped. Good components receive a
second search in the plain subgraph propagating signs. There is no integer
factorization and no expansion of a conjugator.

Let $r=|V|$, $m=|E|$, and let $B\ge1$ bound the bit length of each absolute
edge exponent. A simple tree path has at most $r-1$ edges. A numerator or
denominator formed before cancellation therefore has $O(rB)$ bits.
Fraction reduction cannot increase that bound. Products used in a
fundamental-cycle comparison have the same order of bit size.

\begin{proposition}[Conservative graph bound]\label{prop:graphcost}
The exact inference has an implementation with
\begin{equation}\label{eq:graphcost}
 O\!\left((r+m)\,[rB+\log(r+m+2)]^2\right)
\end{equation}
bit operations, using ordinary quadratic integer arithmetic. A conservative
storage bound is $O(r^2B+mB+(r+m)\log(r+m+2))$ bits. Exponent parsing is
included. Conjugator verification is not included; it is an input obligation.
\end{proposition}
\begin{proof}
There are $O(r+m)$ graph visits and fraction updates. Each update consists
of a bounded number of multiplications, divisions, gcd computations and
comparisons on $O(rB)$-bit integers. Classical quadratic bit bounds suffice.
Dense relabelling and array adjacency avoid requiring an unproved hashing
hypothesis; sorting/comparison-map overhead is absorbed by the displayed
bound. Store at most one fraction and predecessor record per vertex and
one pair of exponents per edge. The sign pass is less expensive.
\end{proof}
The Python reference uses dictionaries, lists and \code{Fraction}. It is
an exact implementation of this arithmetic algorithm, not a claim that
Python dictionary operations or big-integer operations cost one bit step.
The displayed bound is deliberately conservative; improving it is not
necessary for the subsequent polynomial saturation theorem.

\subsection{Independent replay}
A local certificate contains a sorted deleted-vertex list and one cycle
witness per bad full component. A cycle is a list of edge identifiers and
directions, together with the mode \code{modulus} or \code{plain}. It
contains no claimed products that an attacker can inflate or falsify.
The checker validates the walk, its closure and simplicity, recomputes
signed products, checks either $|A|\ne|B|$ or $A\ne B$ with every edge
plain, and independently reconstructs full connectivity by union--find.
It requires the deleted list to equal exactly the union of the justified
components.

Distinct certified components are vertex-disjoint. Taking one simple cycle
per component uses at most $r$ edge occurrences in total (including loops
and two-edge parallel cycles). Thus the combinatorial witness has
$O(r\log(r+m+2))$ bits, apart from the source edge data already supplied.
Balanced product trees check all numeric cycle defects in
$O(\M(rB)\log(r+2))$ bit operations under the usual regular multiplication
cost assumptions; ordinary multiplication gives a polynomial bound without
those assumptions. Source reconstruction, edge authentication and the
connectivity scan must still be charged separately.

The checker intentionally does not rerun fraction discovery or compare its
answer to the producer's answer. It verifies a shorter mathematical
consequence. A proof with no deleted vertices is a harmless local
noncertificate; the source saturation replayer requires a nonempty deletion
in every recorded round.

\subsection{A deterministic finite-field sieve}
Before rational discovery, the optional sieve repeats ratio propagation in
$\F_p^\times$, with $p=65537$. It skips an edge if either exponent vanishes
modulo $p$. A modular inconsistency gives a cycle whose exact integer
products cannot have equal absolute values. The ordinary exact cycle
checker still replays that witness. Every remaining full component then
receives the rational and plain-sign passes.

\begin{proposition}[One-sided sieve with complete fallback]\label{prop:sieve}
The sieve-plus-fallback algorithm returns exactly $\dead$ for every input.
It has no probabilistic equality assumption. Its worst-case bound is not
better than Proposition~\ref{prop:graphcost}, because an inconsistent
integer cycle may be consistent modulo $p$.
\end{proposition}
\begin{proof}
An edge used by the field search has invertible numerator and denominator.
If its discovered cycle had $|A|=|B|$ in $\Z$, their reductions would be
equal in the field, contradicting the discovery. Skipped edges cannot
make a produced cycle false. A missed contradiction is processed by the
exact pass. The sign pass has the same role as before.
\end{proof}

Here the failure is particularly transparent. For $b$ divisible by $32$,
\[
 \frac{2^b+3}{2^b+1}\equiv2\pmod{65537},
\]
since $2^{16}=-1$ and $2^{32}=1$ in the field. The element $2$ has order
$32$. A cycle of length $r$ with that ratio on every edge is inconsistent
over $\Q$, but invisible to this sieve whenever $32\mid r$. A field pass
can therefore add work on a whole natural family of power-of-two-sized
benchmarks. Section~\ref{sec:measurements} retains precisely this control
rather than treating modular agreement as an exact result.

\section{Exposed donors and polynomial deletion-only saturation}
\label{sec:saturation}
\subsection{Literal source extraction}
Let the current presentation have $r$ live generators and fixed relator
slots $R_1,\ldots,R_h$. Write
\[
 N=r+h+\sum_i |R_i|
\]
for its literal size in letters and slots, with generator identifiers
renumbered into $\{1,\ldots,r\}$ for complexity accounting.

The extractor first freely and cyclically reduces each relator and finds
its shortest repeated literal root by the standard prefix-function period
test. A uniform nonempty relator is a single-generator power; in a
knot group it proves that generator trivial and is represented by a
plain self-edge $x^2=x$. Otherwise, for each cyclic signed-run boundary
in a root of length $\ell$, it tests exact decompositions
\begin{equation}\label{eq:pattern}
 x^a\,U\,y^{-b}\,U^{-1},\qquad a,b\ne0,
\end{equation}
with a maximal initial uniform signed run. The length $k=|U|$ determines
the middle run length $\ell-|a|-2k$.

For a fixed rotation, precompute the ends of uniform runs and the longest
common prefix of the candidate $U$ region with the inverse of the whole
rotated word. The latter says exactly which $k$ make the suffix equal
to $U^{-1}$. Each possible $k$ then needs a constant number of indexed
checks. There are $O(\ell)$ rotations and $O(\ell)$ choices per rotation,
so the scan and its number of candidate outputs are $O(\ell^2)$ at the
letter-operation level. Identical edge data are deduplicated. We do not
claim to find every algebraically implied relation of the form
\eqref{eq:pattern}; the claim is exhaustive discovery in this specified
literal exposed class.

A donor proof names its source slot, a root period, a rotation, the head
length and the conjugator length. The separate decoder checks periodicity,
uniform runs and inverse suffixes directly against the current source
relator. It does not use the producer's prefix-function or longest-prefix
routine. A nonminimal but valid repeated-root period would still be sound;
minimality is an extraction optimization, not a trust requirement.

\subsection{All deletion passes have a polynomial bound}
The saturation loop is: extract donors, infer and replay the killed set,
delete its letters from every relator, freely/cyclically reduce, and repeat.
A nonempty pass deletes at least one generator. An empty pass ends the loop.
No generator is introduced, no relator slot is silently removed, and no
word length increases.

\begin{theorem}[Polynomial literal saturation]\label{thm:closure}
On a supplied source-replayed literal presentation of size $N$, complete
saturation of this exposed donor class terminates after at most $r$
nonempty deletion passes and one final unsuccessful pass. Discovery and
independent replay have a conservative bound
\begin{equation}\label{eq:closurecost}
 O\bigl(N^5\log^2(N+2)\bigr)
\end{equation}
in bit operations. A polynomial-size trace suffices. No logarithmic
round-depth hypothesis is needed for this deletion-only closure.
\end{theorem}
\begin{proof}
Deletion followed by reduction never increases the literal size. Every
round has at most $O(N^2)$ deduplicated candidate edges, each exponent has
$O(\log(N+2))$ bits, and there are at most $N$ live generators.
Proposition~\ref{prop:graphcost} is bounded by
$O(N^4\log^2(N+2))$ per pass. Independent donor replay can scan a complete
relator for each of $O(N^2)$ candidates, costing at most
$O(N^3\log(N+2))$ bits per pass; the direct implementation is intentionally
allowed this redundancy. Word deletion and cyclic reduction are smaller.
Even naive worst-case duplicate comparisons remain polynomial and fit the
coarse per-pass bound. Multiply by at most $N+1$ passes. Each donor record
has a bounded number of $O(\log N)$-bit fields, so retaining all round
records is polynomial as well. Soundness is preserved by
Theorem~\ref{thm:sound} at each pass.
\end{proof}

The theorem concerns an already supplied literal presentation. The
braid-to-Artin-word construction can have exponentially long words before
it reaches this stage, and the implementation caps that construction.
Neither the cap nor the polynomial closure theorem converts that frontend
into a complete fast general recognizer.

\subsection{What survives under compressed deletion}
Let a raw straight-line program contain $S$ rules and represent the current
relators, with maximum represented-length bit size $B$. For a deleted
set $D$, the homomorphism $x\mapsto1$ for $x\in D$ and $x\mapsto x$
otherwise is evaluated by one bottom-up pass. It creates at most one
image node per original rule, plus terminals and the empty node. Lengths
cannot increase. Re-evaluating an immutable original grammar using the
cumulative deletion mask avoids counting old allocated nodes as active
new source size.

\begin{proposition}[Conditional compressed closure]\label{prop:compressed}
Suppose a source-checked donor interface can completely inspect these raw
masked grammars in cost $P(S,B,r,h)$, producing at most $m$ edges with
$B_e$-bit exponents and verifying their source witnesses within the same
bound. Then at most $r+1$ such calls, $r$ linear grammar mappings and
$r+1$ applications of Proposition~\ref{prop:graphcost} suffice for deletion
closure. Active raw grammar size and represented-length bit size do not
multiply from one deletion pass to the next.
\end{proposition}
This statement does not provide the assumed donor-discovery interface.
A raw masked grammar can contain cancellations; discovering all the exposed
relations after compressed normalization has a separate cost. The delivered
graph kernel supports binary exponent data, and the literal source wrapper
is implemented, but a complete native WordArena donor collector is not.
The proposition records the useful nonexpansion invariant without hiding
that missing step.

\section{Binary power blocks beyond forest eligibility}
\label{sec:binary}
\subsection{An unconditional gcd-of-minors theorem}
Consider a two-generator group presented solely by plain power equalities
\begin{equation}\label{eq:block}
 H=\left\langle x,y\ \middle|\ x^{a_i}=y^{b_i},\quad
 1\le i\le m\right\rangle,
 \qquad a_i,b_i\in\Z\setminus\{0\}.
\end{equation}
Set $D=\gcd_{i<j}|a_i b_j-a_j b_i|$, taking the gcd of an all-zero family
to be zero.

\begin{theorem}[Exact triviality criterion for two-power blocks]
\label{thm:minors}
The group in Equation~\eqref{eq:block} is trivial if and only if $D=1$.
The implication $D=1\Rightarrow x=y=1$ remains valid when these relations
are a selected subset of a larger presentation, without torsion-freeness
or any knot-group hypothesis.
\end{theorem}
\begin{proof}
For every $i,j$, comparison of equal powers gives
\[
 x^{a_i b_j}=y^{b_i b_j}=x^{a_j b_i},
 \qquad
 y^{b_i a_j}=x^{a_i a_j}=y^{b_j a_i}.
\]
Thus both $x$ and $y$ have every nonzero minor as an annihilating exponent.
A B\'ezout combination of those integers gives $x^D=y^D=1$. In particular,
$D=1$ makes both elements trivial, in any ambient group.

Conversely, the abelianization of $H$ is the quotient of $\Z^2$ by the
lattice spanned by $(a_i,-b_i)$. If this lattice has rank less than two,
that quotient has positive rank and $H$ is nontrivial. At rank two its
index is the gcd of the absolute $2\times2$ minors, namely $D$; this is
the elementary rank-two Smith-normal-form index formula. If $D>1$, the
abelianization is again nontrivial. These alternatives exhaust $D\ne1$.
\end{proof}

A certificate need not contain a unit minor. For the rows
$(2,3),(4,9),(5,4)$, the absolute minors are $6,7,29$; none is one,
but their gcd is one. The delivered checker recomputes only a selected
list of source-row minors and checks their gcd. During production, after
the first nonzero minor, every retained new minor strictly lowers the
positive gcd by a factor of at least two. For $B$-bit coefficients this
needs only $O(B)$ retained minor pairs. Scanning all pairs costs
$O(m^2B^2)$ with ordinary arithmetic; the selected index witness has
$O(B\log(m+2))$ bits and polynomial replay cost. The checker accepts no
unverified claimed determinant or B\'ezout coefficient.

This result does not say that abelianization detects triviality for
arbitrary two-generator presentations. The pure two-power shape is the
entire reason each lattice minor annihilates each individual generator.
Likewise, $D>1$ does not by itself describe the nonabelian structure of $H$.

\begin{proposition}[Collinear rows]\label{prop:collinear}
If all rows in Equation~\eqref{eq:block} are integer multiples
$k_i(a,b)$ of a primitive integer vector $(a,b)$, and
$d=\gcd_i|k_i|$, then the presentation is equivalent to the single relation
$x^{ad}=y^{bd}$.
\end{proposition}
\begin{proof}
Put $u=x^a$, $v=y^b$. The given relations say $u^{k_i}=v^{k_i}$.
Choose $\sum c_i k_i=d$. Multiplying the equalities after taking powers
$c_i$ gives $u^d=v^d$: the left factors are powers of $u$, and the right
factors are powers of $v$, so each side can be combined independently.
Conversely, $d\mid k_i$ implies every original equality from $u^d=v^d$.
No unique-root property has been used.
\end{proof}

\subsection{A binary family that a unit forest cannot contract}
For $M\ge2$ and $k\ge1$, define
\begin{equation}\label{eq:family}
 P_{k,M}=\left\langle t,x_1,y_1,\ldots,x_k,y_k\ \middle|\
 x_i^M=y_i^{M+1},\quad x_i^{M+1}=y_i^{M+2}\ (1\le i\le k)
 \right\rangle.
\end{equation}
Each block has determinant
\[
 M(M+2)-(M+1)^2=-1.
\]
Theorem~\ref{thm:minors} therefore eliminates every $x_i,y_i$ in an
arbitrary group and leaves the free generator $t$. In particular,
$P_{k,M}\cong\Z$ without using external knot provenance.

The word $x^M y^{-(M+1)}$ is not primitive in $F(x,y)$: its one-relator
quotient surjects onto the nonabelian free product $C_M*C_{M+1}$, whereas
killing a primitive word would give a cyclic quotient. The same argument
applies to $x^{M+1}y^{-(M+2)}$. Both have coprime exponent coordinates, so
neither is a proper power. There are no singleton generator occurrences
and no eligible primitive unit-coordinate donors in these displayed
relators. Thus the existing \emph{local donor conditions}, applied to this
unchanged presentation, provide no forest move. A two-edge graph cycle
eliminates each pair immediately.

This is a strict separation of eligible operations, not a lower bound on
Whitehead search, Tietze transformations, or the maintained recognizer.
The family is a supplied presentation of the unknot group; it is not
being passed off as a hard family of input knot diagrams.

When $M=2^b$, the expanded relator list has exactly
$k(4\cdot2^b+4)$ letters, while its two-power encoding has
$O(k(b+\log(k+2)))$ bits. Direct source-row authentication and the
unit-minor checker run in $O(kb^2)$ arithmetic bit cost plus indexing,
without expansion. The graph kernel also remains polynomial in that
binary size. Thus this family gives a fully implemented binary-input
local result, not merely a claim that an unspecified compressed extractor
would be efficient.

\section{A source-bound braid frontend and an independent reference oracle}
\label{sec:source}
\subsection{Reconstructing the entire group}
For a braid with $s$ strands, the closure is a knot exactly when its strand
permutation is a single cycle. The frontend checks strict integer input
and this condition; too few crossings to connect all strands is rejected
before a large permutation allocation. Artin meridian images are then
constructed with an explicit letter allowance. At a positive crossing the
current pair of meridian words $(A,B)$ becomes
\[
 (ABA^{-1},A),
\]
and at a negative crossing it becomes $(B,B^{-1}AB)$. Equating every final
meridian image to its initial meridian gives the full closed-braid group
presentation. This is the usual Artin presentation; see
\cite[Theorem 1.1 and Definition 3.1]{Gemein}.

The producer updates the tuple of current meridians in crossing order.
The independent replayer instead starts with generator words and applies
the corresponding literal substitutions in reverse order. This is the
required reversal for pullback composition, not a discretionary change
of source. All strand relations, including redundant ones, are retained.
Neither replayer invokes the graph producer or the donor extractor.

\subsection{The homological obstruction and a controlled basis change}
\label{sec:homologybarrier}
Initially every Wirtinger or Artin meridian has abelianization charge one.
A generator certified trivial must have charge zero. Therefore an
annihilation-only rule has no direct target among unmodified meridian
generators. This is a structural reason to change basis rather than merely
try more graph cycles on the original meridian graph.

Choose the last meridian as $t$ and make the explicit Nielsen change
\[
 x_s=t,\qquad x_i=y_i t\quad(i<s).
\]
The inverse is $y_i=x_i x_s^{-1}$. All $y_i$ have charge zero and $t$ has
charge one. The substitution at most doubles the literal letter count,
and both source reconstructions independently perform it. Throughout
source replay, any proposal to delete $t$ is rejected. This charge check
is a necessary sanity guard, not the mathematical justification for
other deletions.

After each checked round, all dead letters are removed from every relator
and the surviving words are cyclically reduced. A positive certificate
must end with exactly $t$ alive and every retained relator having exponent
sum zero. Proposition~\ref{prop:terminal} then applies. Failed search
returns \code{INCONCLUSIVE}; a construction allowance failure returns
\code{RESOURCE\_LIMIT}. Neither is a negative knot verdict.

\subsection{A small exact failure of completeness}
The three-braid $\sigma_1^{-1}\sigma_2$ closes to an unknot, by two
successive destabilizations. In the chosen difference basis its reduced
relators, writing $a=y_1$, $b=y_2$, are
\begin{align*}
 R_1&=ba^{-1},\\
 R_2&=t^{-1}b^{-1}atbtb^{-1}t^{-1}a^{-1},\\
 R_3&=t^{-1}b^{-1}atb.
\end{align*}
The exposed scanner finds only the two orientations of $a=b$ from $R_1$.
The graph is consistent, so no deletion occurs. Yet substituting $a=b$
into $R_3$ immediately gives $b=1$, and then $a=1$. The missing move is
an equality contraction, already within the maintained unit-coordinate
machinery. This tiny example disproves completeness of the isolated
annihilation frontend, not completeness of a richer native portfolio.
It also explains why integration should alternate this operation with
existing non-deletion changes rather than replace them.

\subsection{The Khovanov audit is a different implementation}
The reference module resolves all $2^n$ states of a small braid diagram,
constructs circles by literal connectivity, and uses the Frobenius algebra
$\F_2[X]/(X^2)$ with
\[
 m(1,X)=m(X,1)=X,\quad m(X,X)=0,
 \qquad \Delta(1)=1\otimes X+X\otimes1,\quad
 \Delta(X)=X\otimes X.
\]
The marked circle is fixed to label $X$, giving the reduced subcomplex.
The code computes sparse bit-vector boundary ranks and explicitly checks
$d^2=0$. Suppressing gradings does not change total homology dimension:
\[
 \dim\Kh=\sum_h\dim C^h-2\sum_h\operatorname{rank}d_h.
\]
Reduced rational homology has nonzero rank, and the universal coefficient
theorem bounds it by the rank over $\F_2$. Together with
Kronheimer--Mrowka's unknot detection and the unknot calculation, reduced
$\F_2$ rank one is therefore an unknot criterion \cite{KM}. The nonzero
rational rank also follows directly from the normalized Jones Euler
characteristic at $1$.

This oracle is an exponential, capped audit backend. It is not the
maintained Bar--Natan scanning implementation, whose local-complex
approach is a different computational algorithm \cite{BN}. The
standalone command-line option \code{--cube} uses this reference only
after an inconclusive probe and reports a cap failure without a verdict.
Tests of the new code do not constitute execution of the repository's
full maintained test suite.

\section{Finite validation and measured performance}
\label{sec:measurements}
\subsection{Test boundaries}
The shipped standard-library-only suite has \TestCount{} passing test
methods, with no failures, errors or skipped methods in the recorded run.
The methods contain exhaustive and randomized cases, not merely one
assertion apiece. Selected counts are in Table~\ref{tab:tests}.
\begin{table}[htbp]\centering
\begin{tabular}{@{}lr@{}}
\toprule
Check & Count \\
\midrule
Graph/valuation-oracle cases & 3996 \\
Exhaustive short-word donor cases & 5460 \\
Constructed donor cases & 1000 \\
Additional random homology comparisons & 180 \\
Independent Artin reconstructions & 133 \\
Random minor-gcd blocks & 200 \\
Large binary unit-pair checks & 20 \\
\bottomrule
\end{tabular}

\caption{Selected inner-test counts from \code{test\_summary.json}.}
\label{tab:tests}\end{table}
Graph answers are compared against a separate prime-valuation-vector
oracle on small integers, in addition to comparison between exact and
sieved discovery. Every positive graph answer is independently replayed.
The word scanner is compared against a slower all-cut literal oracle.
Other tests cover source rebinding, bad integer types, invalid cycles,
forged cuts, prime-divisible exponents, field collisions, trefoil normal
forms, source construction order, mirror and Markov consistency of the
reference homology, large binary minors and allocation-limit behavior.
These checks support implementation correctness but do not replace the
mathematical proofs or amount to formal verification.

\subsection{Exhaustive short-braid coverage}
\begin{table}[htbp]\centering\small
\begin{tabular}{@{}lrrrr@{}}
\toprule
Source class & Inputs & Unknots & Certified & New$^*$ \\
\midrule
Zero-crossing circle & 1 & 1 & 1 & 0 \\
Two-strand, exhaustive & 170 & 98 & 98 & 0 \\
Three-strand, exhaustive & 2856 & 1712 & 1258 & 0 \\
Four/five-strand, seeded & 200 & 132 & 71 & 0 \\
Total & 3227 & 1943 & 1428 & 0 \\
\bottomrule
\end{tabular}

\caption{Independent small-braid audit. ``New'' means positive coverage beyond
pure-power-only deletion in the same difference basis, not beyond the
maintained recognizer. Inputs are word presentations, not distinct knot types.}
\label{tab:audit}\end{table}
The corpus contains the zero-crossing circle; all signed two-braid words
of lengths $1,3,5,7$; all single-component signed three-braid words of
lengths $2,4,6$; and 200 seeded four/five-strand knot closures with at most
nine crossings. Each row is independently evaluated by the reduced
homology oracle with its square-zero check enabled. All positive group
certificates undergo complete source replay.

Of \AuditCount{} inputs, 1,943 have reduced rank one. The probe certifies
1,428 of them and gives no false positive. It misses 515 unknots. No
source allowance is exhausted. More importantly, the pure-power-only
ablation leaves exactly the same surviving generator set on every input.
Every positive graph certificate seeds its deletion from an exposed
single-generator power. Consequently, this corpus establishes neither
new detection coverage from inconsistent nonseed cycles nor a need for
the new balancedness inference on ordinary small examples. The graph
can propagate a deletion in one batch rather than wait for successive
powers to become exposed, but that is a different claim.

The retained example $\sigma_1^{-1}\sigma_2$ is the first missed unknot;
its full relators were displayed in Section~\ref{sec:source}. It is useful
negative evidence for the frontend design. The largest initial transformed
presentation in the corpus has 290 letters; successful certificates use
at most three deletion rounds. No extrapolation from these sizes to hard
large diagrams is justified.

\subsection{Fixed-prime controls and the follow-up experiment}
All timed graph arms include exact independent certificate replay. Each
initial case has seven shuffled repetitions of exact A/A and sieve A/A
arms, with one warm-up per arm. The nine initial graph cases and twelve
source cases make 588 measured calls and 84 warm-ups. Raw samples,
within-round order, returned work counters, Python/platform metadata,
and source hashes are retained in \code{benchmark.json}.

\begin{table}[H]\centering\small
\begin{tabular}{@{}lrrrrr@{}}
\toprule
Family & $r$ & $b$ & Exact ms & Sieve ms & Paired ratio \\
\midrule
unbalanced & 32 & 64 & 0.622 & 0.932 & 0.605 \\
unbalanced & 128 & 64 & 2.006 & 2.191 & 0.957 \\
unbalanced & 128 & 256 & 7.663 & 7.747 & 1.021 \\
unbalanced & 512 & 64 & 18.164 & 18.941 & 0.957 \\
balanced & 32 & 64 & 0.292 & 0.357 & 0.802 \\
balanced & 128 & 64 & 1.913 & 2.007 & 0.957 \\
balanced & 128 & 256 & 7.500 & 7.373 & 1.005 \\
prime divisible & 128 & 64 & 2.590 & 2.600 & 0.910 \\
residue collision & 128 & 16 & 0.840 & 1.247 & 0.680 \\
\bottomrule
\end{tabular}

\caption{Initial graph families. Exact and sieve columns are medians;
``paired ratio'' is the median of within-round exact/sieve time ratios.
Values below one favor exact-only discovery. All times include replay.}
\label{tab:graphinitial}\end{table}
The nominally unbalanced cycles in Table~\ref{tab:graphinitial} have
lengths divisible by 32, exactly the aliasing condition in
Proposition~\ref{prop:sieve}'s discussion. Their field pass does not
remove the rational work. Apparent small timing gains are not evidence
of successful modular pruning; the recorded rational-edge counters and
A/A variation show why. Balanced, prime-divisible and residue-collision
controls likewise retain the fallback cost.

A separate mechanism-focused follow-up brackets multiples of 32 rather
than replacing the initial observations. It has nine repetitions of
four shuffled arms per case, for 324 measured graph calls and 36 warm-ups.
\begin{table}[H]\centering\small
\begin{tabular}{@{}rrrrrrr@{}}
\toprule
$r$ & $b$ & Exact ms & Sieve ms & Ratio & Exact A/A & Sieve A/A \\
\midrule
31 & 64 & 0.288 & 0.154 & 1.769 & 0.992 & 0.895 \\
32 & 64 & 0.315 & 0.351 & 0.838 & 1.019 & 0.925 \\
33 & 64 & 0.309 & 0.165 & 1.906 & 1.052 & 0.982 \\
127 & 256 & 7.499 & 1.460 & 5.525 & 0.968 & 0.981 \\
128 & 256 & 7.771 & 8.578 & 0.934 & 0.884 & 0.942 \\
129 & 256 & 7.348 & 1.373 & 5.480 & 1.008 & 1.028 \\
511 & 64 & 18.750 & 3.340 & 5.758 & 0.992 & 0.973 \\
512 & 64 & 19.102 & 18.905 & 1.000 & 1.030 & 0.954 \\
513 & 64 & 18.723 & 3.342 & 5.546 & 1.033 & 1.001 \\
\bottomrule
\end{tabular}

\caption{Cycle-length follow-up, recorded in \code{sieve\_followup.json}.
The ratio is paired exact/sieve time including replay. Adjacent multiples
of 32 deliberately retain the exact fallback.}
\label{tab:sieve}\end{table}
The nonmultiples produce genuine field witnesses and avoid the long
rational scan; their gains are several times the A/A fluctuations.
The adjacent multiples do not. These are arithmetic-family measurements,
not a distribution of knot-group presentations and not a uniform
asymptotic improvement. The initial and pre-validation-hardening timing
records remain under \path{data/historical/}; current tables are generated
from the current complete rerun, not from a best-of selection of samples.

\subsection{Binary capacity and full reference-pipeline timings}
\begin{table}[htbp]\centering
\begin{tabular}{@{}rrrrr@{}}
\toprule
Pairs $k$ & $b$ & Coefficient bits & Deleted & Median ms \\
\midrule
1 & 64 & 65 & 2 & 0.041 \\
16 & 256 & 257 & 32 & 0.245 \\
128 & 1024 & 1025 & 256 & 2.791 \\
128 & 16384 & 16385 & 256 & 27.651 \\
\bottomrule
\end{tabular}

\caption{Binary family $P_{k,2^b}$. Each timed call constructs source rows,
checks every unit-minor proof, builds the graph and independently replays
graph deletion. Seven samples are retained per row. There is no expanded
baseline completion time assigned to the large cases.}
\label{tab:capacity}\end{table}
The largest case represents $128(4\cdot2^{16384}+4)$ relator letters using
binary coefficients. It verifies all 256 deleted generators while retaining
one isolated generator. This is a supplied-presentation capacity result.
It proves neither that a hard knot diagram has this form nor that a
producer can expose it cheaply.

\begin{table}[htbp]\centering\small
\begin{tabular}{@{}lrlrrr@{}}
\toprule
Source & $n$ & Unknot & Cube ms & Probe+cube ms & Ratio \\
\midrule
circle & 0 & yes & 0.008 & 0.024 & 0.369 \\
one crossing & 1 & yes & 0.024 & 0.082 & 0.261 \\
trefoil & 3 & no & 0.101 & 0.169 & 0.605 \\
T25 & 5 & no & 1.095 & 1.809 & 0.823 \\
circle three & 2 & yes & 0.075 & 0.164 & 0.421 \\
mixed circle & 4 & yes & 0.476 & 0.162 & 2.936 \\
figure eight & 4 & no & 0.330 & 0.574 & 0.667 \\
T34 & 8 & no & 19.960 & 18.862 & 1.042 \\
circle four & 3 & yes & 0.251 & 0.228 & 1.126 \\
sleeve seven & 7 & yes & 11.032 & 0.363 & 29.305 \\
alternating eight & 8 & no & 16.896 & 19.041 & 0.944 \\
five strand circle & 4 & yes & 1.025 & 0.301 & 3.359 \\
\bottomrule
\end{tabular}

\caption{Twelve fixed small sources: the same exponential reference cube alone
versus the source-replayed probe followed, when needed, by that cube.
These are not comparisons against \code{fastunknot}.}
\label{tab:pipeline}\end{table}
The source timings include positive proof construction and independent
source replay. The reference cube itself already costs only microseconds
on the smallest examples, so a checked group probe can be slower. The
seven-crossing sleeve example benefits by avoiding a full explicit cube,
but contains obvious inverse crossing pairs. Its improvement is not
hard-unknot evidence. The inconclusive cases pay extra work before the
same cube. The A/A records remain available even where they show
substantial submillisecond variation. No broad maintained-recognizer gain
or change of production default is recommended on this evidence.

\section{Integration contract for the maintained recognizer}
\label{sec:integration}
The archive is an additive research package, not an installed modification
of the pinned repository. Its certificate version 1 is local to the
standalone braid wrapper and must not be confused with any maintained
certificate-version number. The intended native handoff has four parts.

First, the maintained replayer reconstructs the caller's original PD and
replays the complete preceding presentation trace. Only that step supplies
classical single-component knot provenance. Second, an exact donor adapter
binds each power-conjugacy edge to a current relator slot and validates its
root and conjugator cuts. The graph verifier must not accept a caller's
unproven \code{plain} flag or a guessed two-power interpretation of a
compressed word. Third, the new graph record gives cycles and the exact
component union to be deleted. Finally, the native checker substitutes
identity into \emph{every} current root and updates the live set before
continuing with existing endpoints.

\paragraph{Policy.}
The small-source ablation suggests keeping this optional. Existing cheap
singletons and unit-coordinate contractions should not be displaced by a
more elaborate graph pass. A promising trigger is the presence of several
verified nonunit power-conjugacy edges in a component that has no already
known pure-power seed. The trigger must be cheap enough not to repeat the
full donor search merely to decide whether to run it. Conversely, once a
useful graph exists, batch propagation can replace repeated identical
connectivity work.

\paragraph{Safe interruption.}
Failure to find a witness means only that this abstraction is exhausted.
Cancellation, node/work allowances and failed compressed cut verification
must preserve the original valid state and return control to existing
fallbacks. Deletion should be committed only after the entire batch is
validated. The theorem does not authorize deleting a relator because it
looked like a donor; it authorizes deleting certified identity generators
and applying that substitution everywhere.

\paragraph{What was not run.}
No remote files were changed, no native WordArena adapter was installed,
and the full maintained recognizer test suite was not executed for this
package. The source-bound braid wrapper, its tests and its independent
homology oracle are new self-contained code. The package includes an
integration checklist rather than an untested production patch.

\section{Further research questions and a route toward the global target}
\label{sec:questions}
\subsection{Compressed donor discovery}
The most immediate algorithmic problem is to find and certify exposed
forms $x^a U y^{-b}U^{-1}$ in a current compressed relator without
enumerating expanded cut positions. Source equality of a proposed cut
and discovery of all useful cuts are distinct tasks. Can candidate cuts
be organized as arithmetic progressions, periodic families or a small
set of grammar-boundary alignments? A useful theorem would bound both
candidate enumeration and rejected-candidate work in the grammar size
and represented-length bit size, and would remain valid after masking
identity generators.

\subsection{Does the cycle rule help on genuinely stalled source states?}
The finite audit found no coverage beyond power deletion. The next
experiment should use actual maintained search states immediately before
an expensive Whitehead/overlap exposure step, not more randomly inflated
circles. For every detected nonunit component, record whether a short
cycle is new information, whether the old native planner already had an
available contraction, and whether full source production plus replay
beats the unchanged baseline. A useful outcome could be positive or
negative: proving rarity would guide dispatch as much as finding gains.

\subsection{Use the labels that the sharpness theorem discards}
Theorem~\ref{thm:sharp} says exactly why the graph alone has reached its
limit. Retain a compressed cyclic conjugator product, prove that it is
trivial or centralizes the base element, or exploit a peripheral marking.
Then some cycles with opposite signed products may become safe plain-type
certificates. Can such extra annotations be verified in polynomial
encoded cost without requiring a general group-word solver? Any proposed
extension must fail on the trefoil inverse-conjugacy example unless its
extra premise genuinely excludes that model.

\subsection{Geometric provenance for stronger signed inference}
A verified absence of generalized torsion of order two permits signed
constraints on all edges. Existing geometric characterizations in
\cite{HMT} provide a theorem-level premise. The algorithmic question is
whether obtaining that premise is cheaper than the recognition work it
saves, and how a native certificate should authenticate it. It is not
legitimate to assume it merely because a diagram appears hyperbolic.

\subsection{Multi-generator unconditional power blocks}
Theorem~\ref{thm:minors} is exact for the two-generator pure-power class.
What is the strongest comparably inexpensive unconditional identity rule
for larger plain power graphs? Unit-defect cycles kill their own cycle
vertices without torsion-freeness, but propagating to an attached vertex
with nonunit exponent needs an additional argument. A generalization
must not confuse a trivial abelianization with a trivial nonabelian group.
Short B\'ezout/minor certificates and provenance-free replay are natural
objects to investigate.

\subsection{Amortized arithmetic and modular scheduling}
The fixed-prime experiment exposes a precise failure mode, not a random
hash accident. Multiple carefully chosen moduli, cheap bit-length tests
on cycle products, or progress-sensitive disabling of unproductive field
passes could improve constants. A deterministic implementation must still
charge all moduli and preserve exact fallback. The more ambitious goal
is a near-linear-in-input-bit-size weighted graph algorithm with comparably
small certificates; exponent factorization should not be introduced as
an uncharged oracle.

\subsection{The remaining quasi-polynomial obligation}
Suppose an input diagram of size $n$ is transformed, with source replay,
into polynomially many presentation states whose total encoded size and
source-checked exposure cost are at most $2^{O((\log n)^c)}$. Assume useful
exposures alternate with the deletion closure above until a decisive
endpoint is reached. Then the displayed local polynomial bounds keep the
whole alternating computation quasi-polynomial. This implication is a
cost-accounting statement, not a proof of those hypotheses.

The useful improvement over a naive round-count argument is that
identity-deletion phases themselves neither introduce generators nor
expand the current literal words, and raw masked grammar evaluation has
the analogous active-size invariant. One need not demand logarithmic
\emph{deletion} depth. The hard missing theorem concerns the availability
and size of decisive exposure transformations, including those needed
on knot diagrams where every graph in the current basis is consistent.
A structural theorem guaranteeing such controlled exposure, rather than
a faster already polynomial cycle check, is what would advance the
unrestricted asymptotic target.

\section{Conclusion}
Power-conjugacy cycles give a valid new operation beyond the acyclic
unit-coordinate forest interface. The exact graph-level boundary is now
explicit: absolute inconsistencies in the full graph and signed
inconsistencies in the plain subgraph generate all universally forced
identities visible to this abstraction. A trefoil countermodel proves
that extending the signed rule without extra hypotheses is impossible.
Independent short certificates and deletion nonexpansion turn this rule
into a polynomial saturation primitive, while binary power blocks provide
fully checked restricted-input examples beyond the old donor eligibility.

The empirical conclusion is narrower. The delivered small-braid frontend
is sound on the audited corpus but incomplete even on a two-crossing
unknot, and the new nonseed cycle mechanism adds no small-corpus coverage.
Modular arithmetic helps on selected bit-complexity families and fails
systematically on adjacent controls. These boundaries identify the next
research task: source-bound, size-controlled exposure of genuinely useful
nonunit cycles in the maintained search, not a claim that the general
quasi-polynomial recognition problem has been completed here.

\appendix
\section{Certificate schema and executable examples}
The graph certificate is deliberately small:
\begin{lstlisting}
{"killed": [1, 2],
 "witnesses": [{"mode": "modulus", "walk": [[0, 1], [1, -1]]}]}
\end{lstlisting}
Its validity depends on the supplied authenticated edges. A direction is
a strict integer $+1$ or $-1$; Booleans are not accepted as integers. A
\code{plain} witness requires all traversed edges to be plain. A
\code{modulus} witness is checked by exact absolute products even when
field arithmetic found it. The deletion list must be sorted, duplicate-free
and exactly the union of the justified connected components.

The standalone source certificate additionally contains the fixed basis
change, every donor list for every round, and the terminal generator.
The original braid is supplied independently to replay. A local graph
proof or two-power minor proof alone is never an unknot certificate.

From the package root:
\begin{lstlisting}
python code/cli.py examples/circle.json --proof /tmp/circle-proof.json
python code/cli.py examples/circle.json --verify /tmp/circle-proof.json
python code/cli.py examples/figure_eight.json --cube
python code/run_tests.py --output /tmp/power-tests.json
python code/audit_braids.py --output /tmp/power-audit.json
python code/benchmark.py --output /tmp/power-benchmark.json
python code/benchmark_sieve.py --output /tmp/power-sieve.json
python verify_manifest.py
sh build.sh
\end{lstlisting}
The last command compiles the article from LaTeX. \code{make\_tables.py}
regenerates numerical tables from the retained JSON without rerunning
measurements. Re-running benchmarks gives new machine-dependent samples;
the data and prose must be reviewed together before replacing the shipped
observations. The tests need Python 3.10 or later and no external Python
packages. Building the paper additionally needs \code{pdflatex} and the
packages in the preamble.

\section{Provenance and reproducibility notes}
The repository read is a pinned-source inspection, not a full checkout
or a native performance baseline. The corresponding source inventory is
in \path{integration/SOURCE_AUDIT.md}. The article's new mathematics and
code were prepared for this delivery. Primary literature is attributed
in the references below. Bibliographic dates should not be confused with
the runtime timestamps: the research-request date is 8 October 2026,
while the retained machine timestamps are UTC on 9 October 2026.

The main raw files are \path{data/test_summary.json},
\path{data/braid_audit.json}, \path{data/benchmark.json}, and
\path{data/sieve_followup.json}. Each timing record includes source hashes;
all measured calls and warm-ups are retained. The historical records are
explicitly separate from the current figures. Archive hashes are supplied
by \path{MANIFEST.sha256}; it intentionally excludes itself and transient
LaTeX/Python build products. The archive does not contain a Lean
formalization or claim machine-checked proofs of the mathematical theorems.

\begin{thebibliography}{9}
\bibitem{HMT}
K. Himeno, K. Motegi and M. Teragaito,
\emph{Generalized torsion, unique root property and Baumslag--Solitar
relation for knot groups}, Hiroshima Mathematical Journal \textbf{53}
(2023), 345--358. The numbered statements cited here refer to
arXiv:2208.00621v1 (2022), especially Lemma 2.3, Example 3.1 and
Corollary 1.8. \url{https://arxiv.org/abs/2208.00621}.

\bibitem{Lackenby}
M. Lackenby, \emph{Incompressible surfaces, hierarchies and unknot
recognition}, arXiv:2607.23350v1 (2026), Section 9 and Proposition 9.1.
\url{https://arxiv.org/abs/2607.23350}.

\bibitem{Gemein}
B. Gemein, \emph{Representations of the singular braid monoid and group
invariants of singular knots}, Topology and its Applications \textbf{114}
(2001), 117--140. The classical Artin statements used here are Theorem 1.1
and Definition 3.1.
\url{https://www.maths.ed.ac.uk/~v1ranick/papers/gemein.pdf}.

\bibitem{Papakyriakopoulos}
C. D. Papakyriakopoulos, \emph{On Dehn's lemma and the asphericity of
knots}, Annals of Mathematics (2) \textbf{66} (1957), 1--26.
The cyclic-group criterion and Loop-Theorem argument are also explicitly
recorded in the inspected repository's group-certificate section.

\bibitem{KM}
P. B. Kronheimer and T. S. Mrowka,
\emph{Khovanov homology is an unknot-detector}, Publications
Math\'ematiques de l'IH\'ES \textbf{113} (2011), 97--208.
\url{https://arxiv.org/abs/1005.4346}.

\bibitem{BN}
D. Bar--Natan, \emph{Fast Khovanov homology computations},
Journal of Knot Theory and Its Ramifications \textbf{16} (2007), 243--255.
\url{https://arxiv.org/abs/math/0606318}.

\bibitem{Repo}
V. Reshetnikov, \emph{ProveIt: Topology/UnknotRecognition}, inspected at
commit \texttt{13d60da8e683ca02b11c6b301c068bda929014e7}.
Relevant files: \path{synthesis/primitive_forest.tex},
\path{synthesis/primitive_projection.tex}, \path{synthesis/group_certificates.tex},
\path{synthesis/report.tex}, \path{synthesis/README.md}, and
\path{reports/45/README.md}.
\url{https://github.com/VladimirReshetnikov/ProveIt}.
\end{thebibliography}
\end{document}
